\documentclass[12pt, a4paper]{article}

\usepackage[utf8]{inputenc}
\usepackage[T1]{fontenc}
\usepackage[spanish, es-tabla]{babel}
\usepackage[landscape, margin=1.2cm]{geometry}
\usepackage{newtxtext, newtxmath}
\usepackage{tikz}
\usetikzlibrary{arrows.meta, positioning, shapes.geometric}
\usepackage{graphicx}
\usepackage{caption}
\usepackage[hidelinks]{hyperref}
\pagestyle{empty}

\begin{document}

\begin{center}
{\bfseries\Large Flujograma para seleccionar la prueba estadística en la contrastación de hipótesis}
\end{center}

\vspace{0.2cm}

\begin{figure}[h!]
\centering
\resizebox{0.98\linewidth}{!}{%
\begin{tikzpicture}[
  node distance=1.05cm and 1.25cm,
  every node/.style={font=\small},
  startstop/.style={
    rectangle,
    rounded corners,
    draw,
    align=center,
    text width=3.3cm,
    minimum height=0.9cm
  },
  process/.style={
    rectangle,
    draw,
    align=center,
    text width=3.4cm,
    minimum height=0.9cm
  },
  decision/.style={
    diamond,
    draw,
    align=center,
    aspect=2.15,
    text width=3.0cm,
    inner sep=1pt
  },
  note/.style={
    rectangle,
    draw,
    align=center,
    text width=4.1cm,
    minimum height=0.75cm
  },
  arrow/.style={-{Stealth[length=2.2mm]}, thick}
]

% Main path
\node[startstop] (start) {Inicio\\Definir hipótesis};
\node[process, right=of start] (desc) {Siempre aplicar\\estadística descriptiva};
\node[decision, right=of desc] (objective) {Objetivo del análisis};

% Descriptive-only branch
\node[process, above=2.0cm of objective] (descriptive) {Un solo grupo\\sin comparación\\Describir resultados};

% Correlation branch
\node[decision, below=2.1cm of objective] (corrnormal) {Correlación\\¿normalidad?};
\node[process, right=of corrnormal] (pearson) {Pearson\\si normal};
\node[process, below=of pearson] (spearman) {Spearman\\si no normal\\u ordinal};

% Comparison branch
\node[decision, right=2.0cm of objective] (scale) {Escala de la\\variable};

% Qualitative branch
\node[decision, above=1.9cm of scale] (qualrel) {Cualitativa\\¿muestras relacionadas?};
\node[process, right=of qualrel] (mcnemar) {McNemar};
\node[process, below=of mcnemar] (chi) {Chi-cuadrada\\Fisher si esperado $<5$};

% Quantitative branch
\node[decision, below=1.9cm of scale] (groups) {Cuantitativa\\número de grupos};
\node[decision, right=1.6cm of groups] (two) {2 grupos\\o 2 momentos};
\node[decision, below=1.35cm of two] (three) {$\geq 3$ grupos\\o mediciones};

% Two-group branch
\node[decision, right=1.55cm of two] (normal2) {¿Normalidad?\\Shapiro-Wilk $n\leq50$\\Kolmogorov-Smirnov $n>50$};
\node[decision, right=1.65cm of normal2] (rel2normal) {Normal\\¿relacionadas?};
\node[decision, below=1.35cm of rel2normal] (rel2nonnormal) {No normal\\¿relacionadas?};

\node[process, right=1.45cm of rel2normal] (tpaired) {t-Student\\relacionada};
\node[process, below=of tpaired] (tind) {t-Student\\independiente};
\node[process, right=1.45cm of rel2nonnormal] (wilcoxon) {Wilcoxon};
\node[process, below=of wilcoxon] (mann) {U de\\Mann-Whitney};

% Three-or-more branch
\node[decision, right=1.55cm of three] (normal3) {¿Normalidad?};
\node[decision, right=1.65cm of normal3] (rel3normal) {Normal\\¿relacionadas?};
\node[decision, below=1.35cm of rel3normal] (rel3nonnormal) {No normal\\¿relacionadas?};

\node[process, right=1.45cm of rel3normal] (anova2) {ANOVA\\2 vías / medidas repetidas};
\node[process, below=of anova2] (anova1) {ANOVA\\1 vía};
\node[process, right=1.45cm of rel3nonnormal] (friedman) {Friedman};
\node[process, below=of friedman] (kruskal) {Kruskal-Wallis};

% Final note
\node[note, below=1.6cm of groups] (justify) {Justificar en el protocolo:\\diseño, mediciones, escala, normalidad y prueba elegida};

% Arrows
\draw[arrow] (start) -- (desc);
\draw[arrow] (desc) -- (objective);
\draw[arrow] (objective) -- node[left] {Describir} (descriptive);
\draw[arrow] (objective) -- node[left] {Correlación} (corrnormal);
\draw[arrow] (objective) -- node[above] {Comparar} (scale);

\draw[arrow] (corrnormal) -- node[above] {Normal} (pearson);
\draw[arrow] (corrnormal) |- node[right] {No normal} (spearman);

\draw[arrow] (scale) -- node[right] {Cualitativa} (qualrel);
\draw[arrow] (scale) -- node[right] {Cuantitativa} (groups);

\draw[arrow] (qualrel) -- node[above] {Sí} (mcnemar);
\draw[arrow] (qualrel) |- node[right] {No} (chi);

\draw[arrow] (groups) -- node[above] {2} (two);
\draw[arrow] (groups) |- node[right] {$\geq3$} (three);

\draw[arrow] (two) -- (normal2);
\draw[arrow] (normal2) -- node[above] {Sí} (rel2normal);
\draw[arrow] (normal2) |- node[right] {No} (rel2nonnormal);

\draw[arrow] (rel2normal) -- node[above] {Sí} (tpaired);
\draw[arrow] (rel2normal) |- node[right] {No} (tind);
\draw[arrow] (rel2nonnormal) -- node[above] {Sí} (wilcoxon);
\draw[arrow] (rel2nonnormal) |- node[right] {No} (mann);

\draw[arrow] (three) -- (normal3);
\draw[arrow] (normal3) -- node[above] {Sí} (rel3normal);
\draw[arrow] (normal3) |- node[right] {No} (rel3nonnormal);

\draw[arrow] (rel3normal) -- node[above] {Sí} (anova2);
\draw[arrow] (rel3normal) |- node[right] {No} (anova1);
\draw[arrow] (rel3nonnormal) -- node[above] {Sí} (friedman);
\draw[arrow] (rel3nonnormal) |- node[right] {No} (kruskal);

\draw[arrow] (groups) -- (justify);

\end{tikzpicture}%
}
\caption*{Fuente: Elaboración propia con base en criterios de diseño del estudio, número de mediciones, escala de variables y supuesto de normalidad.}
\end{figure}

\end{document}
